\documentclass[a4paper,11pt]{article}
\usepackage[utf8]{inputenc}
\usepackage[T1]{fontenc}
\usepackage[french]{babel}
\usepackage{geometry}
\geometry{left=1.5cm,right=1.5cm,top=1.2cm,bottom=1.2cm}

\usepackage{lmodern}
\usepackage{xcolor}
\definecolor{titleblue}{RGB}{0,51,140}

\usepackage{hyperref}
\hypersetup{colorlinks=true, urlcolor=titleblue, linkcolor=titleblue}

\usepackage{titlesec}
\usepackage{enumitem}

\pagestyle{empty}

\titleformat{\section}{\color{titleblue}\Large\bfseries\uppercase}{}{0em}{}[\color{titleblue}\titlerule]
\titlespacing\section{0pt}{8pt}{4pt}
\setlist[itemize]{leftmargin=*,topsep=1pt,itemsep=1pt}

\begin{document}

\begin{center}
    {\Huge\bfseries\color{titleblue} AISSAOUI Ismail}\\[4pt]
    {\Large Ingénieur d'État en Intelligence Artificielle \& Data Science}\\[6pt]
    ismail.aissaoui.pro@gmail.com $|$ +213 660 70 77 96 $|$ Chlef, Algérie \\
    \href{https://linkedin.com/in/aissaoui-ismail-6a77a92a7}{linkedin.com/in/aissaoui-ismail-6a77a92a7} $|$ \href{https://github.com/i-aissaoui}{github.com/i-aissaoui}
\end{center}

\section{RÉSUMÉ PROFESSIONNEL}
Ingénieur diplômé de l'ESI-SBA, expert en architectures distribuées, modèles génératifs et optimisation métaheuristique. Spécialisé dans les systèmes cyber-physiques et l'inférence sur périphériques contraints. Candidat au doctorat focalisé sur l'IA distribuée et la performance énergétique collaborative.

\section{EXPÉRIENCE PROFESSIONNELLE}
\textbf{Stagiaire Ingénieur IA — Jumeau Numérique \& Edge Computing} \hfill Jan 2026 -- Présent \\
\textbf{Sonatrach (Branche Transport par Canalisations)} \\
\begin{itemize}
    \item Conception d'une architecture duale Edge/Cloud pour le monitoring de turbines à gaz.
    \item Optimisation d'inférence sur périphériques contraints (\textbf{0.04ms}) via \textbf{C++ JIT xLSTM}.
    \item Implémentation d'un substitut génératif (\textbf{Mamba-SSM}) pour simulations temps réel.
    \item Intégration de modèles de physique informée (\textbf{PINN}) pour la cohérence thermodynamique.
\end{itemize}

\textbf{Stagiaire en ingénierie réseau} \hfill Sep 2024 \\
\textbf{Algérie Télécom RMC Center}

\section{FORMATION}
\textbf{École Nationale Supérieure d’Informatique (ESI-SBA)} \hfill 2021 -- 2026 \\
\textbf{Diplôme d'Ingénieur d'État \& Master en IA}

\section{COMPÉTENCES TECHNIQUES}
\begin{itemize}
    \item \textbf{IA Distribuée / Cloud}: Federated Learning (conception), Inférence distribuée, Flower.
    \item \textbf{Optimisation}: PSO, Algorithme Génétique (GA), ACO, Optimisation Bayésienne.
    \item \textbf{Deep Learning}: PyTorch, Mamba-SSM, xLSTM, Transformers, GNN, PINN.
    \item \textbf{Architecture}: C++, Python (FastAPI), Docker, Kubernetes, MLflow, CUDA.
\end{itemize}

\section{PROJETS DE RECHERCHE ET INGÉNIERIE}

\textbf{University Scheduling (NP-Complete Optimization)} \hfill 2025 \\
Génération automatique d'emplois du temps via heuristiques hybrides (\textbf{PSO/ACO/GA}).
\textit{Compétence clé : Optimisation dynamique sous contraintes fortes.}

\textbf{Aura — Reinforcement Learning & Distributed Systems} \hfill 2025 \\
Système de trading adaptatif (DQN) avec pipeline de données temps réel distribué.

\textbf{Geo-RAG — Geospatial Knowledge Retrieval} \hfill 2025 \\
Système RAG intégrant des données multi-sources pour le raisonnement spatial.

\textbf{Advanced NLP Optimization} \hfill 2025 \\
Modération multilingue optimisée via PSO et Bayesian Search pour DistilBERT.

\textbf{Vision Industrielle \& Sécurité (YOLOv8)} \hfill 2025 \\
Détection d'équipements de sécurité et tracking d'objets en temps réel sur flux vidéo.

\section{LANGUES}
Arabe (Maternel), Anglais (Avancé), Français (Courant).

\end{document}
